% PREAMBLE
%===============================%
\documentclass[a4paper,11pt]{article}
% Imported packages
\usepackage[a4paper,top=25mm,bottom=25mm,left=25mm,right=25mm]{geometry}
\usepackage[utf8]{inputenc}
\usepackage[spanish,es-tabla]{babel} % es-tabla para que diga "Tabla" en lugar de "Cuadro"
\usepackage{csquotes}
\usepackage{float}
\usepackage{amsmath}
\usepackage{graphicx}
\usepackage{subcaption}
\usepackage{siunitx}
\usepackage{longtable} % Añadido para permitir tablas que ocupen varias páginas

\usepackage{hyperref}
\hypersetup{
    colorlinks=true,
    linkcolor=blue,
    urlcolor=cyan,
    citecolor=magenta,
    pdftitle={Laboratorio de Bomba de Calor},
}

% References
\usepackage[style=numeric,sorting=none]{biblatex}
% \addbibresource{references.bib} % Descomenta esto cuando tengas tu archivo .bib

% Document data
\def\numeroPractica{5}                                    
\def\tituloPractica{Bomba de calor}  
\def\authorName{Andrés Vinuesa Espinosa y José María Martínez Herrada}                                
\def\groupNumber{B.2.4}                                
\def\fechaPractica{25/03/2026}                             
\def\fechaEntrega{15/04/2026}                              

\title{Práctica \numeroPractica: \tituloPractica}
\author{\authorName \\ Grupo \groupNumber}
\date{Sesión de Laboratorio: \fechaPractica \\ Entrega del informe: \fechaEntrega}

%===============================%
% DOCUMENT
%===============================%
\begin{document}

\maketitle
\begin{abstract}
En este trabajo se analiza experimentalmente el comportamiento térmico y la eficiencia de una bomba de calor. A través del registro temporal de las temperaturas en los focos caliente y frío, se estudia la evolución del sistema hacia un estado estacionario, observándose tendencias asintóticas en \SI{50}{\degreeCelsius} y \SI{0}{\degreeCelsius} respectivamente. Se calculó el coeficiente de operación (COP) tanto experimental como ideal (Carnot), evidenciando que el rendimiento real es significativamente inferior al límite teórico debido a irreversibilidades intrínsecas como pérdidas térmicas, fricción y efectos de transferencia de calor finita. El análisis de incertidumbres y la dispersión de los datos confirman la validez del segundo principio de la termodinámica y permiten caracterizar las fluctuaciones microscópicas presentes en la fase de estabilización del dispositivo.
\end{abstract}
\tableofcontents
\vspace{1cm}



\newpage

\section{Resultados}

\subsection{Evolución de la temperatura en función del tiempo}

A continuación se procede a representar la temperatura de ambos focos de nuestra máquina térmica, en función del tiempo, en total se tuvo la máquina
encedida 30 minutos.
\begin{figure}[htbp]
  \centering
   \includegraphics[width=1\textwidth]{t_vs_t.png}
  \caption{Plot de las temperaturas del foco caliente (rojo) y las del foco frío (azul) en función del tiempo. Las incertidumbres que se pueden ver representadas mediante barras de error son las de su respectiva magnitud en cada eje.}
  \label{fig:temperatures}
\end{figure}

Las barras de error representan el error relativo que se ha tenido que tener en cuenta al medir tanto el tiempo como la temperatura; esta resulta muy pequeña en comparación con los datos puesto que el error que se tiene es menor de una parte de la unidad en ambas variables.

En cuanto a la figura, podemos notar que ambas temperaturas tienden asintóticamente a un valor, la del foco caliente a \SI{50}{\degreeCelsius} y la del foco frío a \SI{0}{\degreeCelsius}. Resulta interesante destacar que frente a esta tendencia asintótica en los valores finales encontramos un patrón menos claro en los iniciales, esto es debido a que las fluctuaciones microscópicas del calor, debidas al flujo turbulento que tiene el refrigerante en el régimen transitorio antes de estar en el estado estacionario, provoca que hasta que la máquina se estabilice, la temperatura puede tanto aumentar como disminuir con el tiempo para ambos focos, mientras que en la fase estable, solo aumenta para el caliente y disminuye para el frío, aunque cada vez lo hace de forma más lenta.

\subsection{Eficiencia de la bomba de calor}
Definimos la eficiencia (coeficiente de operación, COP) experimental como la potencia térmica entregada al foco caliente dividida por la potencia eléctrica suministrada. Si $\dot{Q}_{\mathrm{c}}(t)$ es la tasa de calor suministrada al foco caliente y $P(t)$ es la potencia eléctrica, entonces
\begin{equation}
\varepsilon_{\mathrm{exp}}(t)=\frac{\dot{Q}_{\mathrm{c}}(t)}{P(t)}
=\frac{C_{\mathrm{c}}\,\dfrac{dT_{\mathrm{c}}}{dt}(t)}{P(t)}
\label{eq:cop_exp}
\end{equation}
donde $C_{\mathrm{c}}$ es la capacidad térmica efectiva del foco caliente (en unidades de J/K, por ejemplo masa por calor específico más la contribución del recipiente) y $dT_{\mathrm{c}}/dt$ es la derivada temporal de la temperatura medida para el foco caliente (en K/s). Todas las temperaturas deben usarse en kelvin y las potencias en vatios (W).

El COP ideal (Carnot) asociado a los valores instantáneos de temperatura se puede escribir como
\begin{equation}
\varepsilon_{\mathrm{id}}(t)=\frac{T_{\mathrm{c}}(t)}{T_{\mathrm{c}}(t)-T_{\mathrm{f}}(t)}
\label{eq:cop_id}
\end{equation}
con $T_{\mathrm{c}}$ y $T_{\mathrm{f}}$ en kelvin.

\subsection{Eficiencia experimental y eficiencia ideal}
\begin{figure}[htbp]
  \centering
   \includegraphics[width=1\textwidth]{cop_t_log.png}
  \caption{Plot del COP ideal y experimental frente al tiempo, en escala semilogarítmica. Las incertidumbres que se pueden ver representadas mediante barras de error son las de su respectiva magnitud en cada eje.}
  \label{fig:cop}
\end{figure}

Podemos ver que el COP experimental es significativamente menor que el ideal; en primer lugar no podría físicamente ser mayor porque violaría el segundo principio de la termodinámica y el límite impuesto por el ciclo de Carnot. Esta no idealidad se debe a irreversibilidades internas y pérdidas reales en el sistema: pérdidas térmicas por conducción y convección, fugas de calor, pérdidas eléctricas, fricción y efectos de transferencia de calor finita (tiempos finitos), entre otros.
Además resulta interesante ver como el COP experimental, tiende a 0.8, esto nos indica, que en la realidad la máquina térmica deja de ser rentable al alcanzar el régimen estacionario,
ya que consume más energía que la que entrega, esto se debe a las irreversibilidades anteriormente mencionadas. 
Las incertidumbres representan el error relativo que se ha tenido en cuenta a la hora de calcular este coeficiente. En el caso del ideal es notablemente menor ya que solo depende de la temperatura, mientras que en el experimental entra en juego la incertidumbre expandida del trabajo, que a su vez depende de la temperatura y el tiempo. Como estas son valores fijos, son notablemente más grandes en los valores más pequeños, tal como se puede observar en la gráfica.

Asimismo, vemos que en el coeficiente experimental se observa una dispersión de los datos relativamente alta. Esto es debido a las imperfecciones macroscópicas y microscópicas mencionadas con anterioridad; este ruido sigue un patrón aleatorio, pero los \textit{outliers} tienden a compensarse unos con otros, lo cual nos sirve para confirmar que es debido a imperfecciones y no a un fallo sistemático del experimento.

\section{Conclusiones}
Tras el análisis de los datos obtenidos sobre el funcionamiento de la bomba de calor, se extraen las siguientes conclusiones:

Se ha verificado que las temperaturas de ambos focos tienden asintóticamente a valores fijos (\SI{50}{\degreeCelsius} para el foco caliente y \SI{0}{\degreeCelsius} para el frío). La presencia de fluctuaciones en los estados iniciales del experimento se atribuye a la estabilización macroscópica de la máquina y a fenómenos de agitación térmica molecular, que desaparecen una vez alcanzado el régimen de funcionamiento estable.    
La comparación entre el COP experimental ($\varepsilon_{\mathrm{exp}}$) y el ideal ($\varepsilon_{\mathrm{id}}$) confirma las predicciones de la termodinámica clásica. En ningún caso el rendimiento real superó al límite de Carnot, situándose notablemente por debajo de este. Esta discrepancia se justifica por la existencia de procesos irreversibles como la conducción térmica no deseada, fugas de calor al entorno y pérdidas de energía por fricción mecánica y resistencia eléctrica.    
El experimento permite validar con éxito los principios de transferencia de energía y el comportamiento de las máquinas térmicas reales, demostrando que, aunque el ciclo de Carnot es una cota teórica útil, las condiciones de tiempo finito y las imperfecciones de los materiales condicionan la eficiencia de los sistemas de refrigeración y calefacción.

\section{Anexos}
\subsection{Tablas}

\subsubsection{Medidas de tiempo, potencia y temperaturas}

\begin{longtable}{r r r r}
  \caption{Medidas de tiempo, potencia y temperaturas ajustadas a 2 cifras significativas para las incertidumbres según la GUM.}
  \label{tab:raw-data} \\
  \hline
  $t\ (\mathrm{s})$ & $P\ (\mathrm{W})$ & $T_{\mathrm{C}}\ (\si{\celsius})$ & $T_{\mathrm{F}}\ (\si{\celsius})$ \\
  \hline
\endfirsthead
\multicolumn{4}{c}%
{{\bfseries \tablename\ \thetable{} -- Continuación de la página anterior}} \\
  \hline
  $t\ (\mathrm{s})$ & $P\ (\mathrm{W})$ & $T_{\mathrm{C}}\ (\si{\celsius})$ & $T_{\mathrm{F}}\ (\si{\celsius})$ \\
  \hline
\endhead
  \hline \multicolumn{4}{r}{{Continúa en la siguiente página}} \\
\endfoot
  \hline
\endlastfoot
  $0.00 \pm 0.10$ & $0.00 \pm 0.00$ & $23.00 \pm 0.10$ & $17.00 \pm 0.10$ \\
  $30.00 \pm 0.10$ & $1591.20 \pm 0.30$ & $23.70 \pm 0.10$ & $16.00 \pm 0.10$ \\
  $60.00 \pm 0.10$ & $2736.00 \pm 0.60$ & $24.40 \pm 0.10$ & $16.80 \pm 0.10$ \\
  $90.00 \pm 0.10$ & $6093.90 \pm 0.90$ & $25.30 \pm 0.10$ & $16.70 \pm 0.10$ \\
  $120.00 \pm 0.10$ & $7593.6 \pm 1.2$ & $26.30 \pm 0.10$ & $15.60 \pm 0.10$ \\
  $150.00 \pm 0.10$ & $9408.0 \pm 1.5$ & $26.70 \pm 0.10$ & $15.40 \pm 0.10$ \\
  $180.00 \pm 0.10$ & $11289.6 \pm 1.8$ & $28.10 \pm 0.10$ & $14.80 \pm 0.10$ \\
  $210.00 \pm 0.10$ & $12448.8 \pm 2.1$ & $28.60 \pm 0.10$ & $14.00 \pm 0.10$ \\
  $240.00 \pm 0.10$ & $13610.4 \pm 2.4$ & $29.50 \pm 0.10$ & $13.60 \pm 0.10$ \\
  $270.00 \pm 0.10$ & $18635.4 \pm 2.7$ & $30.30 \pm 0.10$ & $13.40 \pm 0.10$ \\
  $300.00 \pm 0.10$ & $21240.0 \pm 3.0$ & $31.00 \pm 0.10$ & $12.10 \pm 0.10$ \\
  $330.00 \pm 0.10$ & $23159.4 \pm 3.3$ & $32.00 \pm 0.10$ & $12.20 \pm 0.10$ \\
  $360.00 \pm 0.10$ & $25912.8 \pm 3.6$ & $32.60 \pm 0.10$ & $11.30 \pm 0.10$ \\
  $390.00 \pm 0.10$ & $29016.0 \pm 3.9$ & $33.20 \pm 0.10$ & $10.50 \pm 0.10$ \\
  $420.00 \pm 0.10$ & $26334.0 \pm 4.2$ & $34.10 \pm 0.10$ & $10.80 \pm 0.10$ \\
  $450.00 \pm 0.10$ & $27904.5 \pm 4.5$ & $35.00 \pm 0.10$ & $9.70 \pm 0.10$ \\
  $480.00 \pm 0.10$ & $36576.0 \pm 4.8$ & $35.90 \pm 0.10$ & $9.30 \pm 0.10$ \\
  $510.00 \pm 0.10$ & $40157.4 \pm 5.1$ & $36.70 \pm 0.10$ & $9.00 \pm 0.10$ \\
  $540.00 \pm 0.10$ & $40500.0 \pm 5.4$ & $37.40 \pm 0.10$ & $8.40 \pm 0.10$ \\
  $570.00 \pm 0.10$ & $36708.0 \pm 5.7$ & $37.80 \pm 0.10$ & $8.00 \pm 0.10$ \\
  $600.00 \pm 0.10$ & $53856.0 \pm 6.0$ & $38.80 \pm 0.10$ & $7.70 \pm 0.10$ \\
  $630.00 \pm 0.10$ & $56964.6 \pm 6.3$ & $39.60 \pm 0.10$ & $7.30 \pm 0.10$ \\
  $660.00 \pm 0.10$ & $52377.6 \pm 6.6$ & $40.10 \pm 0.10$ & $6.90 \pm 0.10$ \\
  $690.00 \pm 0.10$ & $43642.5 \pm 6.9$ & $40.80 \pm 0.10$ & $6.40 \pm 0.10$ \\
  $720.00 \pm 0.10$ & $55778.4 \pm 7.2$ & $41.50 \pm 0.10$ & $6.10 \pm 0.10$ \\
  $750.00 \pm 0.10$ & $58102.5 \pm 7.5$ & $42.10 \pm 0.10$ & $5.90 \pm 0.10$ \\
  $780.00 \pm 0.10$ & $59950.8 \pm 7.8$ & $42.60 \pm 0.10$ & $5.70 \pm 0.10$ \\
  $810.00 \pm 0.10$ & $62750.7 \pm 8.1$ & $43.30 \pm 0.10$ & $5.20 \pm 0.10$ \\
  $840.00 \pm 0.10$ & $65074.8 \pm 8.4$ & $43.70 \pm 0.10$ & $5.00 \pm 0.10$ \\
  $870.00 \pm 0.10$ & $66868.2 \pm 8.7$ & $44.10 \pm 0.10$ & $4.70 \pm 0.10$ \\
  $900.00 \pm 0.10$ & $58464.0 \pm 9.0$ & $44.80 \pm 0.10$ & $4.50 \pm 0.10$ \\
  $930.00 \pm 0.10$ & $72614.4 \pm 9.3$ & $45.30 \pm 0.10$ & $4.30 \pm 0.10$ \\
  $960.00 \pm 0.10$ & $67929.6 \pm 9.6$ & $45.90 \pm 0.10$ & $4.10 \pm 0.10$ \\
  $990.00 \pm 0.10$ & $75438.0 \pm 9.9$ & $46.10 \pm 0.10$ & $3.80 \pm 0.10$ \\
  $1020.00 \pm 0.10$ & $76908 \pm 10$ & $46.80 \pm 0.10$ & $3.60 \pm 0.10$ \\
  $1050.00 \pm 0.10$ & $80703 \pm 11$ & $47.10 \pm 0.10$ & $3.50 \pm 0.10$ \\
  $1080.00 \pm 0.10$ & $84326 \pm 11$ & $47.50 \pm 0.10$ & $3.30 \pm 0.10$ \\
  $1110.00 \pm 0.10$ & $86669 \pm 11$ & $47.90 \pm 0.10$ & $3.10 \pm 0.10$ \\
  $1140.00 \pm 0.10$ & $89011 \pm 11$ & $48.20 \pm 0.10$ & $3.00 \pm 0.10$ \\
  $1170.00 \pm 0.10$ & $89856 \pm 12$ & $48.70 \pm 0.10$ & $2.90 \pm 0.10$ \\
  $1200.00 \pm 0.10$ & $92964 \pm 12$ & $48.80 \pm 0.10$ & $2.70 \pm 0.10$ \\
  $1230.00 \pm 0.10$ & $95288 \pm 12$ & $49.40 \pm 0.10$ & $2.60 \pm 0.10$ \\
  $1260.00 \pm 0.10$ & $97612 \pm 13$ & $49.70 \pm 0.10$ & $2.50 \pm 0.10$ \\
  $1290.00 \pm 0.10$ & $99936 \pm 13$ & $49.90 \pm 0.10$ & $2.40 \pm 0.10$ \\
  $1320.00 \pm 0.10$ & $102260 \pm 13$ & $50.20 \pm 0.10$ & $2.30 \pm 0.10$ \\
  $1350.00 \pm 0.10$ & $104584 \pm 14$ & $50.70 \pm 0.10$ & $2.20 \pm 0.10$ \\
  $1380.00 \pm 0.10$ & $88044 \pm 14$ & $50.90 \pm 0.10$ & $2.10 \pm 0.10$ \\
  $1410.00 \pm 0.10$ & $89958 \pm 14$ & $50.90 \pm 0.10$ & $2.00 \pm 0.10$ \\
  $1440.00 \pm 0.10$ & $111557 \pm 14$ & $51.30 \pm 0.10$ & $2.00 \pm 0.10$ \\
  $1470.00 \pm 0.10$ & $133888 \pm 15$ & $51.50 \pm 0.10$ & $1.80 \pm 0.10$ \\
  $1500.00 \pm 0.10$ & $116205 \pm 15$ & $51.70 \pm 0.10$ & $1.80 \pm 0.10$ \\
  $1530.00 \pm 0.10$ & $118529 \pm 15$ & $52.10 \pm 0.10$ & $1.70 \pm 0.10$ \\
  $1560.00 \pm 0.10$ & $122756 \pm 16$ & $52.20 \pm 0.10$ & $1.70 \pm 0.10$ \\
  $1590.00 \pm 0.10$ & $126182 \pm 16$ & $52.50 \pm 0.10$ & $1.70 \pm 0.10$ \\
  $1620.00 \pm 0.10$ & $126490 \pm 16$ & $52.90 \pm 0.10$ & $1.60 \pm 0.10$ \\
  $1650.00 \pm 0.10$ & $127826 \pm 17$ & $52.80 \pm 0.10$ & $1.60 \pm 0.10$ \\
  $1680.00 \pm 0.10$ & $106848 \pm 17$ & $53.10 \pm 0.10$ & $1.50 \pm 0.10$ \\
  $1710.00 \pm 0.10$ & $160466 \pm 17$ & $53.30 \pm 0.10$ & $1.50 \pm 0.10$ \\
  $1740.00 \pm 0.10$ & $114005 \pm 17$ & $53.60 \pm 0.10$ & $1.40 \pm 0.10$ \\
  $1770.00 \pm 0.10$ & $137122 \pm 18$ & $53.50 \pm 0.10$ & $1.40 \pm 0.10$ \\
  $1800.00 \pm 0.10$ & $142848 \pm 18$ & $54.00 \pm 0.10$ & $1.30 \pm 0.10$ \\
\end{longtable}

\subsubsection{Cálculo de las eficiencias experimental e ideal}

\begin{longtable}{c c c c c c}
\caption{Cálculo de las eficiencias experimental e ideal en función de la diferencia de temperatura.}
\label{tab:T-eficiencias} \\
\hline
$T_C - T_F\ (\si{\degreeCelsius})$ & $Q_F\ (\mathrm{J})$ & $Q_C\ (\mathrm{J})$ & $W\ (\mathrm{J})$ & COP experimental & COP ideal \\
\hline
\endfirsthead
\multicolumn{6}{c}%
{{\bfseries \tablename\ \thetable{} -- Continuación de la página anterior}} \\
\hline
$T_C - T_F\ (\si{\degreeCelsius})$ & $Q_F\ (\mathrm{J})$ & $Q_C\ (\mathrm{J})$ & $W\ (\mathrm{J})$ & COP experimental & COP ideal \\
\hline
\endhead
\hline \multicolumn{6}{r}{{Continúa en la siguiente página}} \\
\endfoot
\hline
\endlastfoot
$6.00 \pm 0.14$ & $0 \pm 0$ & $0 \pm 0$ & $0.00 \pm 0.00$ &  & $49 \pm 1.2$ \\
$7.70 \pm 0.14$ & $1340 \pm 590$ & $2930 \pm 590$ & $1591.20 \pm 0.30$ & $1.84 \pm 0.37$ & $38.55 \pm 0.70$ \\
$7.60 \pm 0.14$ & $3120 \pm 590$ & $5860 \pm 590$ & $2736.00 \pm 0.60$ & $2.14 \pm 0.22$ & $39.15 \pm 0.72$ \\
$8.60 \pm 0.14$ & $3530 \pm 590$ & $9620 \pm 590$ & $6093.90 \pm 0.90$ & $1.579 \pm 0.097$ & $34.70 \pm 0.56$ \\
$10.70 \pm 0.14$ & $6210 \pm 590$ & $13810 \pm 590$ & $7593.6 \pm 1.2$ & $1.818 \pm 0.078$ & $27.99 \pm 0.36$ \\
$11.30 \pm 0.14$ & $6070 \pm 590$ & $15480 \pm 590$ & $9408.0 \pm 1.5$ & $1.645 \pm 0.063$ & $26.54 \pm 0.33$ \\
$13.30 \pm 0.14$ & $10050 \pm 590$ & $21340 \pm 590$ & $11289.6 \pm 1.8$ & $1.890 \pm 0.052$ & $22.65 \pm 0.24$ \\
$14.60 \pm 0.14$ & $10980 \pm 600$ & $23430 \pm 600$ & $12448.8 \pm 2.1$ & $1.882 \pm 0.048$ & $20.67 \pm 0.20$ \\
$15.90 \pm 0.14$ & $13580 \pm 590$ & $27190 \pm 590$ & $13610.4 \pm 2.4$ & $1.998 \pm 0.043$ & $19.04 \pm 0.17$ \\
$16.90 \pm 0.14$ & $11910 \pm 600$ & $30540 \pm 600$ & $18635.4 \pm 2.7$ & $1.639 \pm 0.032$ & $17.96 \pm 0.15$ \\
$18.90 \pm 0.14$ & $12230 \pm 590$ & $33470 \pm 590$ & $21240.0 \pm 3.0$ & $1.576 \pm 0.028$ & $16.09 \pm 0.12$ \\
$19.80 \pm 0.14$ & $14500 \pm 600$ & $37660 \pm 600$ & $23159.4 \pm 3.3$ & $1.626 \pm 0.026$ & $15.41 \pm 0.11$ \\
$21.30 \pm 0.14$ & $14250 \pm 600$ & $40160 \pm 600$ & $25912.8 \pm 3.6$ & $1.550 \pm 0.023$ & $14.354 \pm 0.092$ \\
$22.70 \pm 0.14$ & $13670 \pm 580$ & $42680 \pm 580$ & $29016.0 \pm 3.9$ & $1.471 \pm 0.020$ & $13.496 \pm 0.081$ \\
$23.30 \pm 0.14$ & $20120 \pm 580$ & $46450 \pm 580$ & $26334.0 \pm 4.2$ & $1.764 \pm 0.022$ & $13.187 \pm 0.077$ \\
$25.30 \pm 0.14$ & $22300 \pm 590$ & $50200 \pm 590$ & $27904.5 \pm 4.5$ & $1.799 \pm 0.021$ & $12.180 \pm 0.065$ \\
$26.60 \pm 0.14$ & $17410 \pm 590$ & $53990 \pm 590$ & $36576.0 \pm 4.8$ & $1.476 \pm 0.016$ & $11.618 \pm 0.059$ \\
$27.70 \pm 0.14$ & $17150 \pm 600$ & $57300 \pm 600$ & $40157.4 \pm 5.1$ & $1.427 \pm 0.015$ & $11.186 \pm 0.055$ \\
$29.00 \pm 0.14$ & $19760 \pm 610$ & $60260 \pm 610$ & $40500.0 \pm 5.4$ & $1.488 \pm 0.015$ & $10.709 \pm 0.050$ \\
$29.80 \pm 0.14$ & $25220 \pm 590$ & $61930 \pm 590$ & $36708.0 \pm 5.7$ & $1.687 \pm 0.016$ & $10.435 \pm 0.047$ \\
$31.10 \pm 0.14$ & $12230 \pm 590$ & $66080 \pm 590$ & $53856.0 \pm 6.0$ & $1.227 \pm 0.011$ & $10.031 \pm 0.043$ \\
$32.30 \pm 0.14$ & $12480 \pm 570$ & $69440 \pm 570$ & $56964.6 \pm 6.3$ & $1.219 \pm 0.010$ & $9.683 \pm 0.040$ \\
$33.20 \pm 0.14$ & $19170 \pm 580$ & $71550 \pm 580$ & $52377.6 \pm 6.6$ & $1.366 \pm 0.011$ & $9.435 \pm 0.038$ \\
$34.40 \pm 0.14$ & $30810 \pm 610$ & $74450 \pm 610$ & $43642.5 \pm 6.9$ & $1.706 \pm 0.014$ & $9.126 \pm 0.036$ \\
$35.40 \pm 0.14$ & $21640 \pm 610$ & $77420 \pm 610$ & $55778.4 \pm 7.2$ & $1.388 \pm 0.011$ & $8.888 \pm 0.034$ \\
$36.20 \pm 0.14$ & $21790 \pm 580$ & $79890 \pm 580$ & $58102.5 \pm 7.5$ & $1.375 \pm 0.010$ & $8.709 \pm 0.032$ \\
$36.90 \pm 0.14$ & $22060 \pm 600$ & $82010 \pm 600$ & $59950.8 \pm 7.8$ & $1.368 \pm 0.010$ & $8.557 \pm 0.031$ \\
$38.10 \pm 0.14$ & $22210 \pm 560$ & $84960 \pm 560$ & $62750.7 \pm 8.1$ & $1.3540 \pm 0.0090$ & $8.306 \pm 0.029$ \\
$38.70 \pm 0.14$ & $21540 \pm 590$ & $86610 \pm 590$ & $65074.8 \pm 8.4$ & $1.3310 \pm 0.0090$ & $8.187 \pm 0.028$ \\
$39.40 \pm 0.14$ & $21400 \pm 600$ & $88270 \pm 600$ & $66868.2 \pm 8.7$ & $1.3200 \pm 0.0090$ & $8.052 \pm 0.027$ \\
$40.30 \pm 0.14$ & $32740 \pm 580$ & $91200 \pm 580$ & $58464.0 \pm 9.0$ & $1.560 \pm 0.010$ & $7.890 \pm 0.026$ \\
$41.00 \pm 0.14$ & $20700 \pm 580$ & $93310 \pm 580$ & $72614.4 \pm 9.3$ & $1.2850 \pm 0.0080$ & $7.767 \pm 0.025$ \\
$41.80 \pm 0.14$ & $27850 \pm 610$ & $95780 \pm 610$ & $67929.6 \pm 9.6$ & $1.4100 \pm 0.0090$ & $7.633 \pm 0.024$ \\
$42.30 \pm 0.14$ & $21200 \pm 600$ & $96640 \pm 600$ & $75438.0 \pm 9.9$ & $1.2810 \pm 0.0080$ & $7.547 \pm 0.024$ \\
$43.20 \pm 0.14$ & $22690 \pm 620$ & $99600 \pm 620$ & $76908 \pm 10$ & $1.2950 \pm 0.0080$ & $7.406 \pm 0.023$ \\
$43.60 \pm 0.14$ & $20100 \pm 570$ & $100800 \pm 570$ & $80703 \pm 11$ & $1.2490 \pm 0.0070$ & $7.345 \pm 0.022$ \\
$44.20 \pm 0.14$ & $18210 \pm 590$ & $102540 \pm 590$ & $84326 \pm 11$ & $1.2160 \pm 0.0070$ & $7.255 \pm 0.022$ \\
$44.80 \pm 0.14$ & $17510 \pm 610$ & $104180 \pm 610$ & $86669 \pm 11$ & $1.2020 \pm 0.0070$ & $7.166 \pm 0.021$ \\
$45.20 \pm 0.14$ & $16470 \pm 620$ & $105480 \pm 620$ & $89011 \pm 11$ & $1.1850 \pm 0.0070$ & $7.110 \pm 0.021$ \\
$45.80 \pm 0.14$ & $17700 \pm 630$ & $107560 \pm 630$ & $89856 \pm 12$ & $1.1970 \pm 0.0070$ & $7.027 \pm 0.020$ \\
$46.10 \pm 0.14$ & $14970 \pm 560$ & $107930 \pm 560$ & $92964 \pm 12$ & $1.1610 \pm 0.0060$ & $6.984 \pm 0.020$ \\
$46.80 \pm 0.14$ & $15150 \pm 570$ & $110440 \pm 570$ & $95288 \pm 12$ & $1.1590 \pm 0.0060$ & $6.892 \pm 0.019$ \\
$47.20 \pm 0.14$ & $14060 \pm 590$ & $111670 \pm 590$ & $97612 \pm 13$ & $1.1440 \pm 0.0060$ & $6.840 \pm 0.019$ \\
$47.50 \pm 0.14$ & $12590 \pm 600$ & $112530 \pm 600$ & $99936 \pm 13$ & $1.1260 \pm 0.0060$ & $6.801 \pm 0.019$ \\
$47.90 \pm 0.14$ & $11560 \pm 610$ & $113820 \pm 610$ & $102260 \pm 13$ & $1.1130 \pm 0.0060$ & $6.751 \pm 0.019$ \\
$48.50 \pm 0.14$ & $11300 \pm 630$ & $115880 \pm 630$ & $104584 \pm 14$ & $1.1080 \pm 0.0060$ & $6.677 \pm 0.018$ \\
$48.80 \pm 0.14$ & $28700 \pm 620$ & $116750 \pm 620$ & $88044 \pm 14$ & $1.3260 \pm 0.0070$ & $6.640 \pm 0.018$ \\
$48.90 \pm 0.14$ & $26810 \pm 630$ & $116770 \pm 630$ & $89958 \pm 14$ & $1.2980 \pm 0.0070$ & $6.627 \pm 0.018$ \\
$49.30 \pm 0.14$ & $6800 \pm 560$ & $118360 \pm 560$ & $111557 \pm 14$ & $1.0610 \pm 0.0050$ & $6.581 \pm 0.018$ \\
$49.70 \pm 0.14$ & $-14590 \pm 540$ & $119290 \pm 540$ & $133888 \pm 15$ & $0.8910 \pm 0.0040$ & $6.532 \pm 0.017$ \\
$49.90 \pm 0.14$ & $3830 \pm 580$ & $120040 \pm 580$ & $116205 \pm 15$ & $1.0330 \pm 0.0050$ & $6.510 \pm 0.017$ \\
$50.40 \pm 0.14$ & $3200 \pm 590$ & $121730 \pm 590$ & $118529 \pm 15$ & $1.0270 \pm 0.0050$ & $6.453 \pm 0.017$ \\
$50.50 \pm 0.14$ & $-610 \pm 610$ & $122140 \pm 610$ & $122756 \pm 16$ & $0.9950 \pm 0.0050$ & $6.443 \pm 0.017$ \\
$50.80 \pm 0.14$ & $-2780 \pm 630$ & $123410 \pm 630$ & $126182 \pm 16$ & $0.9780 \pm 0.0050$ & $6.410 \pm 0.017$ \\
$51.30 \pm 0.14$ & $-1390 \pm 630$ & $125100 \pm 630$ & $126490 \pm 16$ & $0.9890 \pm 0.0050$ & $6.356 \pm 0.016$ \\
$51.20 \pm 0.14$ & $-3200 \pm 640$ & $124630 \pm 640$ & $127826 \pm 17$ & $0.9750 \pm 0.0050$ & $6.366 \pm 0.016$ \\
$51.60 \pm 0.14$ & $19130 \pm 640$ & $125970 \pm 640$ & $106848 \pm 17$ & $1.1790 \pm 0.0060$ & $6.323 \pm 0.016$ \\
$51.80 \pm 0.14$ & $-33700 \pm 640$ & $126770 \pm 640$ & $160466 \pm 17$ & $0.7900 \pm 0.0040$ & $6.302 \pm 0.016$ \\
$52.20 \pm 0.14$ & $14020 \pm 570$ & $128030 \pm 570$ & $114005 \pm 17$ & $1.1230 \pm 0.0050$ & $6.260 \pm 0.016$ \\
$52.10 \pm 0.14$ & $-9460 \pm 550$ & $127660 \pm 550$ & $137122 \pm 18$ & $0.9310 \pm 0.0040$ & $6.270 \pm 0.016$ \\
$52.70 \pm 0.14$ & $-13140 \pm 570$ & $129710 \pm 570$ & $142848 \pm 18$ & $0.9080 \pm 0.0040$ & $6.208 \pm 0.015$ \\
\end{longtable}

\begin{longtable}{r c c}
\caption{Evolución del COP experimental e ideal a lo largo del tiempo.}
\label{tab:cop-timeseries} \\
\hline
$t\ (\mathrm{s})$ & COP experimental & COP ideal \\
\hline
\endfirsthead
\multicolumn{3}{c}%
{{\bfseries \tablename\ \thetable{} -- Continuación de la página anterior}} \\
\hline
$t\ (\mathrm{s})$ & COP experimental & COP ideal \\
\hline
\endhead
\hline \multicolumn{3}{r}{{Continúa en la siguiente página}} \\
\endfoot
\hline
\endlastfoot
$0.00 \pm 0.10$ &  & $49 \pm 1.2$ \\
$30.00 \pm 0.10$ & $1.84 \pm 0.37$ & $38.55 \pm 0.70$ \\
$60.00 \pm 0.10$ & $2.14 \pm 0.22$ & $39.15 \pm 0.72$ \\
$90.00 \pm 0.10$ & $1.579 \pm 0.097$ & $34.70 \pm 0.56$ \\
$120.00 \pm 0.10$ & $1.818 \pm 0.078$ & $27.99 \pm 0.36$ \\
$150.00 \pm 0.10$ & $1.645 \pm 0.063$ & $26.54 \pm 0.33$ \\
$180.00 \pm 0.10$ & $1.890 \pm 0.052$ & $22.65 \pm 0.24$ \\
$210.00 \pm 0.10$ & $1.882 \pm 0.048$ & $20.67 \pm 0.20$ \\
$240.00 \pm 0.10$ & $1.998 \pm 0.043$ & $19.04 \pm 0.17$ \\
$270.00 \pm 0.10$ & $1.639 \pm 0.032$ & $17.96 \pm 0.15$ \\
$300.00 \pm 0.10$ & $1.576 \pm 0.028$ & $16.09 \pm 0.12$ \\
$330.00 \pm 0.10$ & $1.626 \pm 0.026$ & $15.41 \pm 0.11$ \\
$360.00 \pm 0.10$ & $1.550 \pm 0.023$ & $14.354 \pm 0.092$ \\
$390.00 \pm 0.10$ & $1.471 \pm 0.020$ & $13.496 \pm 0.081$ \\
$420.00 \pm 0.10$ & $1.764 \pm 0.022$ & $13.187 \pm 0.077$ \\
$450.00 \pm 0.10$ & $1.799 \pm 0.021$ & $12.180 \pm 0.065$ \\
$480.00 \pm 0.10$ & $1.476 \pm 0.016$ & $11.618 \pm 0.059$ \\
$510.00 \pm 0.10$ & $1.427 \pm 0.015$ & $11.186 \pm 0.055$ \\
$540.00 \pm 0.10$ & $1.488 \pm 0.015$ & $10.709 \pm 0.050$ \\
$570.00 \pm 0.10$ & $1.687 \pm 0.016$ & $10.435 \pm 0.047$ \\
$600.00 \pm 0.10$ & $1.227 \pm 0.011$ & $10.031 \pm 0.043$ \\
$630.00 \pm 0.10$ & $1.219 \pm 0.010$ & $9.683 \pm 0.040$ \\
$660.00 \pm 0.10$ & $1.366 \pm 0.011$ & $9.435 \pm 0.038$ \\
$690.00 \pm 0.10$ & $1.706 \pm 0.014$ & $9.126 \pm 0.036$ \\
$720.00 \pm 0.10$ & $1.388 \pm 0.011$ & $8.888 \pm 0.034$ \\
$750.00 \pm 0.10$ & $1.375 \pm 0.010$ & $8.709 \pm 0.032$ \\
$780.00 \pm 0.10$ & $1.368 \pm 0.010$ & $8.557 \pm 0.031$ \\
$810.00 \pm 0.10$ & $1.3540 \pm 0.0090$ & $8.306 \pm 0.029$ \\
$840.00 \pm 0.10$ & $1.3310 \pm 0.0090$ & $8.187 \pm 0.028$ \\
$870.00 \pm 0.10$ & $1.3200 \pm 0.0090$ & $8.052 \pm 0.027$ \\
$900.00 \pm 0.10$ & $1.560 \pm 0.010$ & $7.890 \pm 0.026$ \\
$930.00 \pm 0.10$ & $1.2850 \pm 0.0080$ & $7.767 \pm 0.025$ \\
$960.00 \pm 0.10$ & $1.4100 \pm 0.0090$ & $7.633 \pm 0.024$ \\
$990.00 \pm 0.10$ & $1.2810 \pm 0.0080$ & $7.547 \pm 0.024$ \\
$1020.00 \pm 0.10$ & $1.2950 \pm 0.0080$ & $7.406 \pm 0.023$ \\
$1050.00 \pm 0.10$ & $1.2490 \pm 0.0070$ & $7.345 \pm 0.022$ \\
$1080.00 \pm 0.10$ & $1.2160 \pm 0.0070$ & $7.255 \pm 0.022$ \\
$1110.00 \pm 0.10$ & $1.2020 \pm 0.0070$ & $7.166 \pm 0.021$ \\
$1140.00 \pm 0.10$ & $1.1850 \pm 0.0070$ & $7.110 \pm 0.021$ \\
$1170.00 \pm 0.10$ & $1.1970 \pm 0.0070$ & $7.027 \pm 0.020$ \\
$1200.00 \pm 0.10$ & $1.1610 \pm 0.0060$ & $6.984 \pm 0.020$ \\
$1230.00 \pm 0.10$ & $1.1590 \pm 0.0060$ & $6.892 \pm 0.019$ \\
$1260.00 \pm 0.10$ & $1.1440 \pm 0.0060$ & $6.840 \pm 0.019$ \\
$1290.00 \pm 0.10$ & $1.1260 \pm 0.0060$ & $6.801 \pm 0.019$ \\
$1320.00 \pm 0.10$ & $1.1130 \pm 0.0060$ & $6.751 \pm 0.019$ \\
$1350.00 \pm 0.10$ & $1.1080 \pm 0.0060$ & $6.677 \pm 0.018$ \\
$1380.00 \pm 0.10$ & $1.3260 \pm 0.0070$ & $6.640 \pm 0.018$ \\
$1410.00 \pm 0.10$ & $1.2980 \pm 0.0070$ & $6.627 \pm 0.018$ \\
$1440.00 \pm 0.10$ & $1.0610 \pm 0.0050$ & $6.581 \pm 0.018$ \\
$1470.00 \pm 0.10$ & $0.8910 \pm 0.0040$ & $6.532 \pm 0.017$ \\
$1500.00 \pm 0.10$ & $1.0330 \pm 0.0050$ & $6.510 \pm 0.017$ \\
$1530.00 \pm 0.10$ & $1.0270 \pm 0.0050$ & $6.453 \pm 0.017$ \\
$1560.00 \pm 0.10$ & $0.9950 \pm 0.0050$ & $6.443 \pm 0.017$ \\
$1590.00 \pm 0.10$ & $0.9780 \pm 0.0050$ & $6.410 \pm 0.017$ \\
$1620.00 \pm 0.10$ & $0.9890 \pm 0.0050$ & $6.356 \pm 0.016$ \\
$1650.00 \pm 0.10$ & $0.9750 \pm 0.0050$ & $6.366 \pm 0.016$ \\
$1680.00 \pm 0.10$ & $1.1790 \pm 0.0060$ & $6.323 \pm 0.016$ \\
$1710.00 \pm 0.10$ & $0.7900 \pm 0.0040$ & $6.302 \pm 0.016$ \\
$1740.00 \pm 0.10$ & $1.1230 \pm 0.0050$ & $6.260 \pm 0.016$ \\
$1770.00 \pm 0.10$ & $0.9310 \pm 0.0040$ & $6.270 \pm 0.016$ \\
$1800.00 \pm 0.10$ & $0.9080 \pm 0.0040$ & $6.208 \pm 0.015$ \\
\end{longtable}

\subsection{Incertidumbres}
En primer lugar, se han analizado las incertidumbres de las variables medidas de forma directa. Apoyándonos en las especificaciones de los instrumentos, se ha establecido una incertidumbre absoluta directa para la temperatura de $u_T = \SI{0.10}{\kelvin}$ (aplicable tanto a $T_{\mathrm{c}}$ como a $T_{\mathrm{f}}$) y para el tiempo de $u_t = \SI{0.10}{\second}$.

A continuación, para las medidas indirectas, la incertidumbre expandida absoluta se calcula empleando la ley de propagación de incertidumbres mediante el formalismo de las derivadas parciales. En su forma general, para una magnitud $y = f(x_1, x_2, \dots, x_n)$, la incertidumbre absoluta $u_y$ viene dada por:
\begin{equation}
u_y = \sqrt{\sum_{i=1}^{n} \left( \frac{\partial f}{\partial x_i} \right)^2 u_{x_i}^2}
\label{eq:unc_general_base}
\end{equation}

\subsubsection{Propagación de incertidumbre para la Potencia y el Trabajo}
Si se asume que el trabajo eléctrico $W$ consumido a lo largo del tiempo depende de la potencia y el tiempo transcurrido ($W = P \cdot t$), su incertidumbre absoluta se calcula como:
\begin{equation}
u_W = \sqrt{\left(\frac{\partial W}{\partial P}\right)^2 u_P^2 + \left(\frac{\partial W}{\partial t}\right)^2 u_t^2} = \sqrt{t^2 u_P^2 + P^2 u_t^2}
\end{equation}

De forma análoga, suponiendo que la potencia $P$ se obtiene como medida indirecta de la tensión $V$ y la corriente $I$ ($P = V \cdot I$), su incertidumbre absoluta correspondiente será:
\begin{equation}
u_P = \sqrt{\left(\frac{\partial P}{\partial V}\right)^2 u_V^2 + \left(\frac{\partial P}{\partial I}\right)^2 u_I^2} = \sqrt{I^2 u_V^2 + V^2 u_I^2}
\end{equation}

\subsubsection{Propagación de incertidumbre para el COP}
Para el COP experimental aplicamos la propagación por derivadas parciales:
\begin{equation}
u_{\varepsilon_{\mathrm{exp}}}=\sqrt{\left(\frac{\partial\varepsilon}{\partial C_{\mathrm{c}}}\right)^2 u_C^2
 +\left(\frac{\partial\varepsilon}{\partial (dT_{\mathrm{c}}/dt)}\right)^2 u_{dT_{\mathrm{c}}/dt}^2
 +\left(\frac{\partial\varepsilon}{\partial P}\right)^2 u_P^2}
\label{eq:unc_prop_general_cop}
\end{equation}
con las derivadas parciales
\begin{align}
\frac{\partial\varepsilon}{\partial C_{\mathrm{c}}}&=\frac{dT_{\mathrm{c}}/dt}{P}, \\
\frac{\partial\varepsilon}{\partial (dT_{\mathrm{c}}/dt)}&=\frac{C_{\mathrm{c}}}{P}, \\
\frac{\partial\varepsilon}{\partial P}&=-\frac{C_{\mathrm{c}}\,dT_{\mathrm{c}}/dt}{P^2}=-\frac{\varepsilon_{\mathrm{exp}}}{P}.
\end{align}
Sustituyendo en (\ref{eq:unc_prop_general_cop}) se obtiene la expresión final de incertidumbre absoluta para la eficiencia experimental:
\begin{equation}
u_{\varepsilon_{\mathrm{exp}}}=\sqrt{\left(\frac{dT_{\mathrm{c}}/dt}{P}\right)^2 u_{C_\mathrm{c}}^2
 +\left(\frac{C_{\mathrm{c}}}{P}\right)^2 u_{dT_{\mathrm{c}}/dt}^2 +\left(\frac{\varepsilon_{\mathrm{exp}}}{P}\right)^2 u_P^2 }.
\label{eq:unc_exp}
\end{equation}

Para el COP ideal (ecuación \ref{eq:cop_id}) la propagación absoluta viene dada por
\begin{equation}
u_{\varepsilon_{\mathrm{id}}}=\sqrt{\left(\frac{\partial\varepsilon_{\mathrm{id}}}{\partial T_{\mathrm{c}}}\right)^2 u_{T_{\mathrm{c}}}^2
 +\left(\frac{\partial\varepsilon_{\mathrm{id}}}{\partial T_{\mathrm{f}}}\right)^2 u_{T_{\mathrm{f}}}^2}
\end{equation}
con
\begin{align}
\frac{\partial\varepsilon_{\mathrm{id}}}{\partial T_{\mathrm{c}}}&=-\frac{T_{\mathrm{f}}}{(T_{\mathrm{c}}-T_{\mathrm{f}})^2},\\
\frac{\partial\varepsilon_{\mathrm{id}}}{\partial T_{\mathrm{f}}}&=\frac{T_{\mathrm{c}}}{(T_{\mathrm{c}}-T_{\mathrm{f}})^2}.
\end{align}
Por tanto, la incertidumbre del COP ideal queda expresada de la siguiente manera:
\begin{equation}
u_{\varepsilon_{\mathrm{id}}}=\sqrt{\frac{T_{\mathrm{f}}^2}{(T_{\mathrm{c}}-T_{\mathrm{f}})^4}u_{T_{\mathrm{c}}}^2
 +\frac{T_{\mathrm{c}}^2}{(T_{\mathrm{c}}-T_{\mathrm{f}})^4}u_{T_{\mathrm{f}}}^2 }.
\label{eq:unc_id}
\end{equation}

% Descomenta la siguiente línea cuando vayas a imprimir tu bibliografía con el archivo .bib
% \printbibliography

\end{document}